% Generated draft: requires mathematical and source review.
\documentclass[11pt]{article}
\usepackage[margin=25mm]{geometry}
\usepackage{fontspec}
\setmainfont{Latin Modern Roman}
\usepackage{amsmath,amssymb,mathtools,unicode-math}
\setmathfont{Latin Modern Math}
\usepackage{enumitem,needspace,xurl,hyperref}
\setlength{\parindent}{0pt}
\setlength{\parskip}{5pt}
\setlength{\emergencystretch}{3em}
\newcommand{\N}{\mathbb{N}}
\newcommand{\Z}{\mathbb{Z}}
\newcommand{\Q}{\mathbb{Q}}
\newcommand{\R}{\mathbb{R}}
\newcommand{\C}{\mathbb{C}}
\newcommand{\F}{\mathbb{F}}
\newcommand{\A}{\mathbb{A}}
\newcommand{\PP}{\mathbb{P}}
\newcommand{\E}{\mathbb{E}}
\newcommand{\eps}{\varepsilon}
\newcommand{\dd}{\,\mathrm{d}}
\begin{document}
\section*{UnsolvedMath ID: 700885}
\textbf{Source problem code:} catalogue entry
\subsection{Definitions and mathematical statement}
The supplied source asks for a structural description of two classes of locally compact soluble groups. The following formulation uses the usual conventions that locally compact groups are Hausdorff and that solubility refers to the underlying abstract group. These conventions are not explicitly specified in the supplied record and require confirmation against the source.

Let \(G\) be a Hausdorff locally compact topological group. Define its abstract derived series by
\[
G^{(0)}=G,\qquad G^{(j+1)}=[G^{(j)},G^{(j)}],
\]
where \([H,H]\) is the subgroup generated by all commutators \(x^{-1}y^{-1}xy\) with \(x,y\in H\), without taking topological closure. Call \(G\) soluble if \(G^{(r)}=\{1\}\) for some integer \(r\geq 0\).

Write
\[
\mathcal{S}(G)=\{H\leq G: H\text{ is closed in }G\text{ and is noncompact in its subspace topology}\}.
\]
The \emph{maximum condition for closed noncompact subgroups} means that every ascending sequence
\[
H_1\leq H_2\leq\cdots,\qquad H_n\in\mathcal{S}(G),
\]
is eventually constant: there exists \(N\geq 1\) such that \(H_n=H_N\) for every \(n\geq N\). The \emph{minimum condition for closed noncompact subgroups} means the corresponding statement for every descending sequence
\[
H_1\geq H_2\geq\cdots,\qquad H_n\in\mathcal{S}(G).
\]
These conditions concern all closed noncompact subgroups, without requiring them to be proper or normal.

\textbf{Problem (V. M. Poletskikh, 1992).} Give structural characterizations, separately, of
\begin{enumerate}
\item the locally compact soluble groups satisfying the maximum condition for closed noncompact subgroups;
\item the locally compact soluble groups satisfying the minimum condition for closed noncompact subgroups.
\end{enumerate}
No noncompactness assumption on \(G\) itself is stated. In particular, compact groups satisfy both chain conditions vacuously. The parenthetical wording in the source is interpreted as two separate classification questions. If the source uses a different separation convention or a notion of solubility involving closures of derived subgroups, the corresponding classification targets must be adjusted; the supplied record does not resolve these conventions.
\subsection{Short English statement}
Describe, separately, the structure of locally compact soluble groups in which every ascending, respectively descending, chain of closed noncompact subgroups eventually stabilizes. Compact ambient groups are included. The draft explicitly uses Hausdorffness and abstract solubility, conventions that require source confirmation.
\subsection{Sources}
\begin{itemize}
\item Local UnsolvedMath registry, permanent ID \(700885\), source problem code \texttt{catalogue entry}, title \emph{Kourovka Notebook Problem 12.67}. The supplied registry labels the problem \texttt{open}; this is a supplied status, not an independently verified mathematical or literature assessment.
\item \emph{The Kourovka Notebook (Unsolved Problems in Group Theory)}, 12th issue (1992), Problem 12.67; proposed by V. M. Poletskikh. This attribution and bibliographic information are supplied by the record.
\item The record also identifies the 21st issue edition (2026), arXiv identifier \texttt{1401.0300v48}. Its collection description names E. I. Khukhro and V. D. Mazurov as editors of the collection covering Issues 1--21 (1965--2026). No external source was independently consulted for this draft.
\end{itemize}
\end{document}
